% [原文 p.95]
\chapter{锥定理}

第 1 章中，我们证明了非奇异射影簇的锥定理，并指出该证明在奇异簇上并不奏效。在极小模型纲领中，某些类型的奇点无法避免，因此准备好对 $(X,\Delta)$ 的锥定理极为重要。这一证明在技术上、历史上都错综复杂，经数人之手才成为现在这个样子。主要贡献见 [Kaw84a, Rei83c, Sho85]。

§3.1 叙述证明的 4 个主要阶段，并说明其背后的基本思想。4 个部分有一个共同之处，称为无基点方法（basepoint-freeness method）。这一方法在无基点定理的证明中最清晰地显现出来。因此，逻辑上应属证明第 2 步的无基点定理，其证明放在 §3.2 中给出。

无基点方法还在许多不同情形下得到应用。其中若干见 [Laz96] 与 [Kol97, Sec.5] 的解说。

其余 3 个阶段在 §3.3 至 §3.5 中处理。其中，有理性定理的证明最为繁复。

§3.6 说明无基点定理与锥定理的相对版本。

这 3 个结果一旦建立，就能精确叙述对数极小模型纲领。这将在 §3.7 中进行。在二维情形，纲领中不会出现翻转。因此，这一情形可以完全处理。

% [原文 p.96]
§3.8 研究对的极小模型。可以看到，这一概念并不是非奇异簇的极小模型的简单推广（定义 2.13）。其定义在定义 3.50 中给出，基本性质则在定理 3.52 中叙述。

\section{锥定理证明引论}
\takeaway{先把消没定理 2.64 改写为定理 3.1，再陈述无基点定理、非消没定理、有理性定理与锥定理；贯穿四者的工具是同一个无基点方法。}

§1.3 中我们已对非奇异簇证明了锥定理。这里先从陈述若干导向一般情形锥定理证明的定理开始。该证明基于一套颇为不同的思想。即使应用于非奇异情形，它也能证明以往方法无法得到的结论，即“极射线总可以收缩”。另一方面，关于生成极射线的曲线，这一证明几乎不提供任何信息；而且该论证只在特征 0 下才起作用。在继续之前，先对消没定理 2.64 稍作改写。

\begin{theorem}{3.1}
设 $Y$ 是非奇异复射影簇，$\sum d_iD_i$ 是（写成不同素除子之和的）$\mathbb{Q}$-除子，$L$ 是线丛（或 Cartier 除子）。设 $D:=L+\sum d_iD_i$ 是 nef 且大的，且 $\sum D_i$ 是 snc 除子。此时，
\[
H^i(Y,\mathcal{O}_Y(K_Y+\lceil D\rceil))=0,\qquad \forall i>0.
\]
\end{theorem}
% [Tutor Note: 定理 3.1 即消没定理 2.64 的改写；$\lceil D\rceil$ 是对 Q-除子的各系数逐项取上整所得的整系数除子，一般不是 D 本身，故不能把 D 直接代入上同调。]

\par\medskip\noindent\textbf{3.2}\quad 陈述 4 个基本定理，最后一个便是锥定理。这 4 个定理的证明，从历史上看相当错综复杂。对非奇异三维簇，[Mor82] 证明了特殊情形。关于三维簇的最初一般结果始于 [Kaw84b]，并由 [Ben83] 与 [Rei83c] 最终完成。非消没定理由 [Sho85] 得到。锥定理始于 [Kaw84a]，完成于 [Kol84]。详细的处理以及向相对情形的推广，请参阅 [KMM87]。\hfill$\square$\par\medskip

\begin{theorem*}[{无基点定理}]{3.3}
设 $(X,\Delta)$ 是 $\Delta$ 有效的固有 klt 对。设 $D$ 是 nef 的 Cartier 除子，且对适当的 $a>0$，$aD-K_X-\Delta$ 是 nef 且大的。此时，对所有的 $b\gg 0$，$|bD|$ 无基点。
\end{theorem*}

% [原文 p.97]
\begin{theorem*}[{非消没定理}]{3.4}
设 $X$ 是固有簇，$D$ 是 nef 的 Cartier 除子，$G$ 是 $\mathbb{Q}$-除子。假设：
\begin{enumerate}[label=(\arabic*),leftmargin=2.4em,itemsep=1pt,topsep=2pt,parsep=0pt]
\item 对适当的 $a>0$，$aD+G-K_X$ 是 $\mathbb{Q}$-Cartier、nef 且大的，
\item $(X,-G)$ 是 klt 的。
\end{enumerate}
此时，对所有的 $m\gg 0$，
\[
H^0(X,mD+\lceil G\rceil)\neq 0.
\]
\end{theorem*}

\begin{theorem*}[{有理性定理}]{3.5}
设 $(X,\Delta)$ 是固有 klt 对，$\Delta$ 有效，且 $K_X+\Delta$ 不是 nef 的。设 $a(X)>0$ 是使 $a(X)\cdot(K_X+\Delta)$ 为 Cartier 的整数。设 $H$ 是 nef 且大的除子，定义
\[
r=r(H):=\max\{t\in\mathbb{R}:H+t(K_X+\Delta)\ \text{是 nef 的}\}.
\]
此时，$r$ 是有理数；在其最简分数表示 $r=u/v$（$u,v\in\mathbb{Z}$，$v>0$）中，
\[
0<v\leq a(X)\cdot(\dim X+1).
\]
\end{theorem*}

\par\medskip\noindent\textbf{补充 3.6}（[KMM94a]）\quad 记号同上。此时存在极射线 $R$，同时满足 $R\cdot(K_X+\Delta)<0$ 与 $R\cdot(H+r(K_X+\Delta))=0$。\hfill$\square$\par\medskip
% [Tutor Note: 有理性定理对每个丰富除子 $H$ 给出临界值 $r(H)$；补充 3.6 说明临界处恰出现 $(K_X+\Delta)$-负的极射线，这正是把锥定理化为抽象凸几何定理 3.15 的桥梁。]

\begin{theorem*}[{锥定理}]{3.7}
设 $(X,\Delta)$ 是 $\Delta$ 有效的射影 klt 对。此时：
\begin{enumerate}[label=(\arabic*),leftmargin=2.4em,itemsep=3pt,topsep=2pt,parsep=0pt]
\item 存在至多可数个有理曲线 $C_j\subset X$，满足 $0<-(K_X+\Delta)\cdot C_j\leq 2\dim X$，且
\[
\overline{NE}(X)=\overline{NE}(X)_{(K_X+\Delta)\geq 0}+\sum\mathbb{R}_{\geq 0}[C_j].
\]
\item 对任意的 $\epsilon>0$ 与丰富的 $\mathbb{Q}$-除子 $H$，
\[
\overline{NE}(X)=\overline{NE}(X)_{(K_X+\Delta+\epsilon H)\geq 0}
+\sum_{\text{有限和}}\mathbb{R}_{\geq 0}[C_j].
\]
\item 设 $F\subset\overline{NE}(X)$ 是 $(K_X+\Delta)$-负的极锥面。此时，在满足 $(\mathrm{cont}_F)_*\mathcal{O}_X=\mathcal{O}_Z$ 的到射影簇的态射 $\mathrm{cont}_F:X\to Z$ 之中，使得对任意不可约曲线 $C\subset X$ 都有“$\mathrm{cont}_F(C)$ 为一点 $\Leftrightarrow[C]\in F$”的态射恰有一个存在。
% [原文 p.98]
$\mathrm{cont}_F$ 称为 $F$ 的收缩（contraction）。
\item 设极锥面 $F$ 与态射 $\mathrm{cont}_F:X\to Z$ 如 (3)。设 $L$ 是 $X$ 上的线丛，且对任何满足 $[C]\in F$ 的曲线 $C$ 都有 $(L\cdot C)=0$。此时存在 $Z$ 上的线丛 $L_Z$，使得 $L\simeq\mathrm{cont}_F^*L_Z$。
\end{enumerate}
\end{theorem*}
% [Tutor Note: 原书把凸锥的面（face）称为「锥面」，把由极射线张成的面称为「端锥面」，字面近于“极锥面”；书中从第 1 章起统一作「极锥面」（extremal face），本条从统一译名，即 $\overline{NE}(X)$ 的面。]

\par\medskip\noindent\textbf{注}\quad 断言 (3) 常被称为收缩定理。\par\medskip

\par\medskip\noindent\textbf{3.8}\quad 这些定理的证明的逻辑顺序如下。
\[
\text{非消没定理}\ \Rightarrow\ \text{无基点定理}\ \Rightarrow\ \text{有理性定理}\ \Rightarrow\ \text{锥定理}.
\]
但是，为了使证明易于理解，我们最先给出无基点定理的证明。因为这一证明最清楚地表达了其中的思想。\hfill$\square$\par\medskip

\par\medskip\noindent\textbf{3.9}（基本方针）\quad 与非消没定理、有理性定理同样，无基点定理的证明的主要思想如下。\par
为简单起见，设 $X$ 非奇异且 $\Delta=0$。设 $M$ 是丰富的 Cartier 除子，$F\subset X$ 是不可约除子。有正合列
\[
0\to\mathcal{O}_X(K_X+M)\to\mathcal{O}_X(K_X+M+F)\to\mathcal{O}_F(K_F+M|_F)\to 0,
\]
由定理 2.64，它给出满射
\[
H^0(X,\mathcal{O}_X(K_X+M+F))\to H^0(F,\mathcal{O}_F(K_F+M|_F)).
\]
注意：即使 $M$ 不是丰富的，只要消没定理可适用于 $K_X+M$，这个映射仍然是满射。因此，若能写成 $bD-K_X=M+F$，则 $bD|_F-K_F=M|_F$，上述满射不外是
\[
H^0(X,\mathcal{O}_X(bD))\to H^0(F,\mathcal{O}_F(bD|_F)),
\]
可以期待用关于维数的归纳法得到截面。一般情形下并非如此。
% [Tutor Note: 3.9 的关键是把 $bD-K_X$ 拆成 nef 且大的一部分再加一个不可约除子 $F$，使消没定理能作用于 $K_X+M$，从而用维数归纳法构造截面。]

但是，在由固有双有理态射 $f:Y\to X$ 拉回后讨论，就有了少许调整的余地。对使 $\lceil A\rceil$ 成为有效 $f$-例外除子的 $Y$ 上的 $\mathbb{Q}$-除子 $A$，有 $H^0(X,\mathcal{O}_X(bD))=H^0(Y,\mathcal{O}_Y(bf^*D+\lceil A\rceil))$。因此，若有满射
\[
H^0(X,\mathcal{O}_X(bD+\lceil A\rceil))\to H^0(F,\mathcal{O}_F(bD|_F+\lceil A\rceil|_F)),
\]
就能与归纳法结合起来。
% [Tutor Note: 此处记号容易混淆：左端印作 $H^0(X,\cO_X(bD+\lceil A\rceil))$，而 $\lceil A\rceil$ 是 $Y$ 上的除子，前一行同型的式子写作 $H^0(Y,\cO_Y(bf^*D+\lceil A\rceil))$；阅读时按 $f$ 的推前来理解，正文照录。]

% [原文 p.99]
因此，把 $M$ 取成使 $bf^*D-K_Y\equiv M+F-\lceil A\rceil$，且使消没定理可适用于 $K_Y+M$；若能做到，就能使用上面的正合列。为应用消没定理，需要
\[
bf^*D-K_Y\equiv(\text{nef 且大的 }\mathbb{Q}\text{-除子})+F-A.
\]
怎样才能找到 $b$ 与 $F$ 呢？必须巧妙地利用给定除子的线性组合。首先，$aD-K_X$ 与 $D$ 可以自由使用。其次，若 $|mD|$ 不是无基点的，则可以使用它分解成无基点的部分 $|L|$ 与固定部分 $\sum r_jF_j$：
\[
f^*|mD|=|L|+\sum r_jF_j.
\]
记 $K_Y=f^*K_X+\sum a_jF_j$。（在和 $\sum a_jF_j$ 中，$a_j>0$ 的只能是例外除子 $F_j$；而和 $\sum r_jF_j$ 中还包含非例外除子。对此必须留意。）
于是，
\[
mf^*D\equiv L+\sum r_jF_j,\qquad
af^*D-K_Y\equiv f^*(aD-K_X)-\sum a_jF_j.
\]
将第 1 式乘以 $c$ 后与第 2 式相加，再加上 $(b-cm-a)f^*D$，便得到下式：
\[
bf^*D-K_Y=(b-cm-a)f^*D+cL+f^*(aD-K_X)-\sum(a_j-cr_j)F_j.
\]
若 $b\geq cm+a$ 且 $c\geq 0$，则右边最初 3 项之和是 nef 且大的。因此，常数 $c$ 必须选得使等式
\[
\sum(a_j-cr_j)F_j=A-F
\]
中 $\lceil A\rceil$ 为有效 $f$-例外除子、$F$ 不可约。此外，还有限制分数系数除子的问题 3.10，因此 $F$ 不能是 $A$ 的成分。
从而，$c$ 要选得使 $\min\{a_j-cr_j\}=-1$。此时，$F$ 等于使 $a_j-cr_j=-1$ 的那些 $F_j$ 的并。这样的 $F_j$ 也许有多个。为了排除这种可能性，用系数非常小的线性组合 $\sum p_jF_j$ 把所有东西都扰动一下。这虽是小技巧，却颇为巧妙。\hfill$\square$\par\medskip

% [原文 p.100]
\par\medskip\noindent\textbf{3.10}\quad 进行以上讨论时，必须考虑把 $\mathbb{Q}$-除子及其取上整限制到非奇异簇 $Y$ 的非奇异素除子 $F$ 上。实际讨论中，只需把形如 $D=L+\sum d_iD_i$ 的除子限制到与任何 $D_i$ 都不同的 $F$，或者限制到与某个 $d_j$ 为整数的 $D_j$ 相一致的 $F$。后一情形下，在限制之前先把 $d_jD_j$ 吸收进 $L$。无论哪种情形，我们都只处理剩余 $D_i$ 之和与 $F$ 相交成简单正规交叉除子的状况。此时，取上整的操作与限制操作可交换。\hfill$\square$\par\medskip
% [Tutor Note: 3.10 限定所考虑的除子形状，使「取上整」与「限制到 F」可交换；这一技术前提在全章反复使用，不可略过。]

\section{无基点定理}
\takeaway{先用对数奇点消解把 $aD-(K_X+\Delta)$ 化为丰富，借非消没定理得到截面；再假设基点集非空，限制到不可约除子 $F$ 上用同一套论证导出矛盾。}

下面的证明几乎原封不动地采用 [Rei83c] 的论述。

\par\medskip\noindent\textbf{步骤 1}\quad 在这一步中，证明对所有的 $m\gg 0$，$|mD|\neq\varnothing$。使用命题 2.61，构造对数奇点消解 $f:Y\to X$，使其满足以下条件：
\begin{enumerate}[label=(\arabic*),leftmargin=2.4em,itemsep=1pt,topsep=2pt,parsep=0pt]
\item $K_Y\equiv f^*(K_X+\Delta)+\sum a_jF_j$，且对所有的 $j$ 有 $a_j>-1$；
\item 对适当的 $a>0$ 与 $0<p_j\ll 1$，$f^*(aD-(K_X+\Delta))-\sum p_jF_j$ 是丰富的。（$F_j$ 不必是 $f$-例外的。）
\end{enumerate}
在 $Y$ 上得到如下等式：
\begin{align*}
f^*(aD-(K_X+\Delta))-\sum p_jF_j
&=af^*D+\sum(a_j-p_j)F_j-\Bigl(f^*(K_X+\Delta)+\sum a_jF_j\Bigr)\\
&\equiv af^*D+G-K_Y.
\end{align*}
其中 $G=\sum(a_j-p_j)F_j$。由假设，$\lceil G\rceil$ 是有效 $f$-例外除子（因 $\Delta$ 有效，$a_j>0$ 只可能在 $F_j$ 为 $f$-例外时出现），$af^*D+G-K_Y$ 是丰富的，并且 $H^0(Y,mf^*D+\lceil G\rceil)=H^0(X,mD)$ 成立。于是应用非消没定理，得到
\[
H^0(X,mD)>0\qquad(\forall m\gg 0).
\]

\par\medskip\noindent\textbf{步骤 2}\quad 对 $s\in\mathbb{Z}_{>0}$，记 $B(s)$ 为 $|sD|$ 的（赋以约化结构的）基点集。对任意正整数 $u>v$，$B(s^u)\subset B(s^v)$ 是显然的。由 Noether 归纳法，（递减）列 $B(s^u)$ 会稳定，将其记为 $B_s$。注意：要么存在 $s>1$ 使 $B_s\neq\varnothing$，要么可以选取两个互素的自然数 $s,s'>1$ 使 $B_s=B_{s'}=\varnothing$。后一情形下，取 $u,v$ 使 $B(s^u)=B(s'^v)=\varnothing$；再利用充分大的自然数可写成 $s^u$ 与 $s'^v$ 的非负整数系数线性组合这一事实，便得到对所有 $m\gg 0$，$|mD|$ 无基点。因此，只需假设某个 $B_s$ 非空并导出矛盾。取 $m=s^u$ 使 $B_s=B(m)$，并假设该集合非空。
% [Tutor Note: 步骤 2 的 Noether 归纳法只用到 $B(s^u)$ 的单调性；若两个互素指数的基点集都空，则由非负系数的线性组合可知所有充分大的 m 都无基点。]

从非消没定理得到的线性系出发，再把奇点消解 $Y$ 进一步爆破，则新的对数奇点消解 $f:Y\to X$ 满足步骤 1 的条件，并且对适当的 $m>0$，表示
\[
f^*|mD|=|L|\ (\text{移动部分})+\sum r_jF_j\ (\text{固定部分})
\]
中出现的 $|L|$ 可假定无基点。因此，$\bigcup\{f(F_j):r_j>0\}$ 是 $|mD|$ 的基点集（记为 $\mathrm{Bs}\,|mD|$）。注意 $f^{-1}\mathrm{Bs}\,|mD|=\mathrm{Bs}\,|mf^*D|$。只要找到满足 $r_j>0$ 的某个 $F_j$，使得对所有的 $b\gg 0$，$F_j$ 都不含于 $|bf^*D|$ 的基点集中，就能如所预期得到矛盾。

% [原文 p.101]
\par\medskip\noindent\textbf{步骤 3}\quad 对满足 $b>cm+a$ 的整数 $b>0$ 与有理数 $c>0$，如下定义除子：
\begin{align*}
N(b,c)&:=bf^*D-K_Y+\sum(-cr_j+a_j-p_j)F_j\\
&\equiv(b-cm-a)f^*D &&\text{（nef）}\\
&\quad +c\Bigl(mf^*D-\sum r_jF_j\Bigr) &&\text{（自由因子，即无基点的）}\\
&\quad +f^*(aD-K_X-\Delta)-\sum p_jF_j &&\text{（丰富的）}.
\end{align*}
按定义，若 $b\geq cm+a$，则 $N(b,c)$ 是丰富的。在这一假设下，由消没定理得 $H^1(Y,\lceil N(b,c)\rceil+K_Y)=0$。其中
\[
\lceil N(b,c)\rceil=bf^*D+\sum\lceil-cr_j+a_j-p_j\rceil F_j-K_Y.
\]

\par\medskip\noindent\textbf{步骤 4}\quad $c$ 与 $p_j$ 可选得使对某个 $F=F_j$，$\sum(-cr_j+a_j-p_j)F_j=A-F$ 成立。其中 $\lceil A\rceil$ 有效，且 $A$ 不以 $F$ 为成分。实际上，先选 $c>0$ 使下式成立：
\[
\min_j(-cr_j+a_j-p_j)=-1.
\]
可以假定达到左边最小值的 $j$ 是唯一的。（若有多个，则预先稍许扰动 $p_j$，即可使 $j$ 唯一。）这个 $j$ 满足 $r_j>0$ 以及
\[
\lceil N(b,c)\rceil+K_Y=bf^*D+\lceil A\rceil-F.
\]
由步骤 3，若 $b\geq cm+a$，则下面的是满射：
\[
H^0(Y,bf^*D+\lceil A\rceil)\to H^0(F,(bf^*D+\lceil A\rceil)|_F).
\]

% [原文 p.102]
\par\medskip\noindent\textbf{注}\quad 若 $F_j$ 出现在 $\lceil A\rceil$ 中，则 $a_j>0$，从而 $F_j$ 是 $f$-例外的。因此 $\lceil A\rceil$ 是 $f$-例外的。\par\medskip

\par\medskip\noindent\textbf{步骤 5}\quad 注意下述等式：
\[
\lceil N(b,c)\rceil|_F=(bf^*D+A-F-K_Y)|_F=(bf^*D+A)|_F-K_F.
\]
由此，把非消没定理应用于 $F$，得到
\[
H^0(F,(bf^*D+\lceil A\rceil)|_F)\neq 0.
\]
因此，$H^0(Y,bf^*D+\lceil A\rceil)$ 有在 $F$ 上不为 0 的元素。由于 $\lceil A\rceil$ 是 $f$-例外的且有效，
\[
H^0(Y,bf^*D+\lceil A\rceil)=H^0(Y,bf^*D)=H^0(X,bD)
\]
成立。由此，如 3.9 中所见，若 $b\gg 0$，则 $f(F)$ 不含于 $|bD|$ 的基点集中。至此，无基点定理的证明完成。

无基点定理最重要的应用之一，是关于标准环的有限生成性。

\begin{theorem*}{3.11}
设 $(X,\Delta)$ 是 $\Delta$ 有效的固有 klt 对。设 $K_X+\Delta$ 是 nef 且大的。此时，由
\[
\bigoplus_{m=0}^{\infty}H^0(X,\mathcal{O}_X(mK_X+\lfloor m\Delta\rfloor))
\]
定义的标准环（canonical ring）在 $\mathbb{C}$ 上有限生成。
\end{theorem*}
% [Tutor Note: 标准环用 $\lfloor m\Delta\rfloor$ 而非 $\lceil m\Delta\rceil$，因为只有向下取整才满足次可加性；这与定理 3.1、3.4 中的上整形成对照。]

\par\medskip\noindent\textbf{注}\quad 标准环实际上是环，因为 $\lfloor m_1\Delta\rfloor+\lfloor m_2\Delta\rfloor\leq\lfloor(m_1+m_2)\Delta\rfloor$。若换成 $\lceil m\Delta\rceil$，不等式反向，因此用 $\lceil m\Delta\rceil$ 得不到环。\par\medskip

\begin{proof}
由定理 3.3，存在 $r>0$，使得 $r\Delta$ 是（整系数）除子且 $\mathcal{O}_X(rK_X+r\Delta)$ 由全局截面生成。这些截面定义态射 $f:X\to Z$，并且存在 $Z$ 上的丰富可逆层 $L$，满足 $f^*L=\mathcal{O}_X(rK_X+r\Delta)$。记 $G_m=f_*\mathcal{O}_X(mK_X+\lfloor m\Delta\rfloor)$。则
\[
\bigoplus_{m=0}^{\infty}H^0(X,\mathcal{O}_X(mK_X+\lfloor m\Delta\rfloor))
=\bigoplus_{m=0}^{\infty}H^0(Z,G_m).
\]
$G_m$ 是凝聚层，且由投影公式 $G_{m+r}\simeq G_m\otimes L$。由于 $L$ 是丰富的，$R=\bigoplus_{m=0}^{\infty}H^0(Z,L^m)$ 是 $\mathbb{C}$ 上的有限生成环，且对所有的 $0\leq j<r$，$\bigoplus_{m=0}^{\infty}H^0(Z,G_{j+rm})$ 是有限生成 $R$-模。于是
\[
\bigoplus_{m=0}^{\infty}H^0(Z,G_m)
=\bigoplus_{j=0}^{r-1}\Bigl(\bigoplus_{m=0}^{\infty}H^0(Z,G_{j+rm})\Bigr)
\]
是 $\mathbb{C}$ 上的有限生成环。
\end{proof}

% [原文 p.103]
上面所述的结果是下述猜想的特殊情形。

\begin{conjecture*}[丰度猜想]{3.12}
设 $(X,\Delta)$ 是 $\Delta$ 有效的固有 lc 对。此时：
\begin{enumerate}[label=(\arabic*),leftmargin=2.4em,itemsep=2pt,topsep=2pt,parsep=0pt]
\item $\bigoplus_{m=0}^{\infty}H^0(X,\mathcal{O}_X(mK_X+\lfloor m\Delta\rfloor))$ 是有限生成环。
\item 若 $K_X+\Delta$ 是 nef 的，则对适当的 $m>0$，$|m(K_X+\Delta)|$ 无基点。
\end{enumerate}
\end{conjecture*}

\par\medskip\noindent\textbf{3.13}\quad 如定理 3.11 中所示，若 $K_X+\Delta$ 是 nef 的，则 (2) 蕴含 (1)。一般地，极小模型纲领把 (1) 归结为 (2)。因此，常常只有 (2) 被称为丰度猜想。\par
即使对曲面，这也不是显然的结果。该猜想在三维已知成立。三维中 $\Delta=0$ 的情形是宫冈（Miyaoka）与川又一系列结果的归结。关于证明的简化与参考文献，请参阅 [K+92]。一般的三维对数标准情形已在 [KMM94a] 中证明。\par
关于更高维的情形，所知甚少。\hfill$\square$\par\medskip
% [Tutor Note: 丰度猜想在三维已成立，二维亦非平凡；曲面情形需先建立极小模型理论，故不能由一般维数论证直接得到。]

\section{锥定理}
\takeaway{把锥定理化归为有理性定理与抽象定理 3.15：丰富除子给出 $\overline{NE}(X)$ 的各个面，逐面下降即得极射线的可数分解。}

本节介绍的锥定理的证明，是通过 J. Kollár、松木谦二、森重文、T. Luo 之间的讨论得到的。

% [原文 p.104]
\par\medskip\noindent\textbf{3.14}（思路说明）\quad 说明有理性定理如何被用来获得关于曲线锥的信息。\par
设 $H$ 是 $X$ 上的丰富除子，按有理性定理取 $r=r(H)$。由于 $r$ 是有理数，对适当的 $m>0$，$m(H+r(K_X+\Delta))$ 是 Cartier 的。注意 $m(H+r(K_X+\Delta))$ 是 nef 的但不是丰富的。因此，$\overline{NE}(X)\cap(H+r(K_X+\Delta))^{\perp}$ 是 $\overline{NE}(X)$ 的一个锥面。从各种丰富除子出发，就得到 $\overline{NE}(X)$ 的各种锥面。锥定理的证明不外是这一考察的形式归结点。准确地说，锥定理可以立即由有理性定理与下面的抽象结果得到。\hfill$\square$\par\medskip
% [Tutor Note: 3.14 中 $H+r(K_X+\Delta)$ 只是 nef 而非丰富，故其零集与 $\overline{NE}(X)$ 的交是一个非平凡的面；变动 H 即得不同的面。]

\begin{theorem*}{3.15}
设 $N_{\mathbb{Z}}$ 是有限秩自由 $\mathbb{Z}$-模，$N_{\mathbb{R}}=N_{\mathbb{Z}}\otimes_{\mathbb{Z}}\mathbb{R}$ 是到 $\mathbb{R}$ 的系数扩张。设 $\overline{NE}\subset N_{\mathbb{R}}$ 是不含任何直线的闭凸锥。设 $K$ 是对偶 $\mathbb{Q}$-向量空间 $N_{\mathbb{Q}}^*$ 的元素，且对适当的 $C\in\overline{NE}$ 有 $(K\cdot C)<0$。假设存在适当的 $a(K)\in\mathbb{Z}_{>0}$，使得对任意的 $H\in N_{\mathbb{Z}}^*$（在 $\overline{NE}-\{0\}$ 上 $H>0$），
\[
r:=\max\{t\in\mathbb{R}:\ \overline{NE}\ \text{上}\ H+tK\geq 0\}
\]
都是形如 $u/a(K)$（$u\in\mathbb{Z}$）的有理数。此时，下述等式成立：
\[
\overline{NE}=\overline{NE}_{K\geq 0}+\sum\mathbb{R}_{\geq 0}[\xi_i].
\]
其中 $\xi_i\in N_{\mathbb{Z}}$ 满足 $(\xi_i\cdot K)<0$，且 $\mathbb{R}_{\geq 0}[\xi_i]$ 在半空间 $K<0$ 上不聚积。
\end{theorem*}
% [Tutor Note: 定理 3.15 是纯凸几何命题；证明中把 K 与非零元 ξ 配对，用有理性定理把系数限制为整数，从而得到极射线的离散性。]

\begin{proof}[定理 3.15、补充 3.6 与锥定理 3.7 的证明]
不妨假设 $K:=K_X+\Delta$ 不是 nef 的。

\par\medskip\noindent\textbf{步骤 1}\quad 对 nef 的除子类 $L$，记 $F_L=L^{\perp}\cap\overline{NE}$。若 $L$ 不是丰富的，则由定理 1.18 得 $F_L\neq\{0\}$。假设 $F_L\not\subset\overline{NE}_{K\geq 0}$，并设 $H$ 是任意的丰富 Cartier 除子。对 $n\in\mathbb{Z}_{>0}$，定义
\[
r_L(n,H):=\max\Bigl\{t\in\mathbb{R}:nL+H+\frac{t}{a(K)}K\ \text{是 nef 的}\Bigr\}.
\]
由假设或者有理性定理，$r_L(n,H)$ 是（非负的）整数；又因 $L$ 是 nef 的，$r_L(n,H)$ 是 $n$ 的非减函数。于是，当 $n\geq n_0$ 时 $r_L(n,H)$ 收敛到固定的 $r_L(H)$。（这是因为，取 $\xi\in F_L\setminus\overline{NE}_{K\geq 0}$，则
\[
r_L(n,H)\leq a(K)\frac{(H\cdot\xi)}{-(K\cdot\xi)},
\]
故可知 $r_L(n,H)$ 是有界非减函数。）
\end{proof}
